% SGA 3, Expose VI_A, Loop 1 English reconstruction.
% Scope: paragraph 4.3 in full; paragraph 4.4 is excluded.
% Authority: Polo--Gille Exp6A-13oct24.pdf, component-local pages 27--28
% (combined-reader pages 330--331; printed pages 319--320).
% OCR and the Reinhold English Markdown are comparison witnesses only.
% Loop 2 replaces the four exact 180-dpi source crops with native
% tikz-cd reconstructions controlled against the same PDF coordinates:
% local p. 27, x=490,y=1235,w=520,h=145;
% local p. 28, x=470,y=525,w=550,h=80;
% local p. 28, x=470,y=695,w=550,h=135;
% local p. 28, x=420,y=920,w=650,h=78;
% all source renders are 1488x2105 pixels.

\medskip
\phantomsection\label{par:sga3-VIA-4.3}%
\noindent\textbf{4.3.---}
由前面的讨论可知, 由 \(u\) 定义的、以 \(G\) 为底的群胚 \(G_*\),
是群胚 \(\mathscr S(G^\alpha_*)\) 的直和;
这里这些群胚由 \(G_*\) 在 \(G\) 的各个形如 \(\mathscr S(G^\alpha)\)
的开闭部分上诱导. 因此, \(G_*\) 的余核是各群胚
\(\mathscr S(G^\alpha_*)\) 的余核的直和,
于是我们只需分别研究这些余核.

先考虑 \(G_*\) 在 \(\mathscr S(G^0)\) 上诱导的群胚
\(\mathscr S(G^0_*)\). 显然, \(\mathscr S(G^0_*)\)
就是由同态
\[
  F\longrightarrow\mathscr S(G^0)
\]
定义的、以 \(\mathscr S(G^0)\) 为底的群胚,
其中这个同态由 \(u\) 诱导 (\S~\hyperref[par:sga3-VIA-3.1]{3.1}).
所以我们要证明其存在的余核等同于 \(F\backslash\mathscr S(G^0)\).
另一方面, 考虑 \(\mathscr S(G^0_*)\) 在 \(G^0\) 上诱导的群胚
\begin{center}
  \phantomsection\label{diag:sga3-VIA-4.3-G0-groupoid}%
  \begin{tikzcd}[column sep=large]
    G^0_2
      \arrow[r,shift left=1.4ex,"\ell'_2"]
      \arrow[r,"\ell'_1"]
      \arrow[r,shift right=1.4ex,"\ell'_0"'] &
    G^0_1
      \arrow[r,shift left=0.8ex,"\ell_1"]
      \arrow[r,shift right=0.8ex,"\ell_0"'] &
    G^0_0=G^0
  \end{tikzcd}
\end{center}
依照第~V~篇 \S~3.b 中的显式构造,
该处记为 \(Y_0\times_{X_0}X_1\) 的对象正是 \(F\times_A G^0\).
因此, \(G^0_1\) 是 \(G^0\) 在由 \(\lambda\) 诱导的态射
\[
  F\times_A G^0\longrightarrow\mathscr S(G^0)
\]
下的逆像.

我们断言, 这个逆像就是 \(F_0\times_A G^0\),
其中 \(F_0\) 表示 \(G^0\) 在 \(u\) 下的逆像.
事实上, 若 \(F^\beta\) 是 \(F_0\) 的一个连通分支,
则由 \SGARefLink{lem:sga3-VIA-2.1.2}{2.1.2},
\(F^\beta\times_A G^0\) 连通,
且 \(\lambda(F^\beta\times_A G^0)\) 包含于 \(G^0\).%
\begingroup
\renewcommand{\thefootnote}{\fnsymbol{footnote}}%
\footnote[1]{英译者注: 法文底本此处印为 \(\lambda(F\times_A G^0)\).
所述断言针对的是刚引入的连通分支 \(F^\beta\),
紧接着的逆向论证也证实, 此处意在限制到 \(F^\beta\times_A G^0\).}%
\endgroup
反过来, 若 \(F^\beta\) 是 \(F\) 的一个不包含于 \(F_0\) 的连通分支,
则 \(F^\beta\times_A G^0\) 的像仍然连通, 且包含 \(u(F^\beta)\);
若 \(u(F^\beta)\) 不包含于 \(G^0\),
则 \(\lambda(F^\beta\times_A G^0)\) 与 \(G^0\) 不相交.

因此, \(G_*\) 在 \(G^0\) 上诱导的群胚 \(G^0_*\)
就是由 \(u\) 诱导的同态
\[
  F_0\longrightarrow G^0
\]
所定义的、以 \(G^0\) 为底的群胚.
由于 \(G^0\), 因而 \(F_0\), 都在 \(A\) 上有限型,
\hyperref[sec:sga3-VIA-3]{第~3~节}表明 \(G^0_*\) 有余核,
即 \(F_0\backslash G^0\).%
\begingroup
\renewcommand{\thefootnote}{\fnsymbol{footnote}}%
\footnote[2]{英译者注: 法文底本写的是“由\hyperref[sec:sga3-VIA-4]{第 4 段}”.
此处引用的有限型商的结果在\hyperref[sec:sga3-VIA-3]{第~3~节};
\hyperref[sec:sga3-VIA-4]{第~4~节}是当前正在进行的一般情形的构造.}%
\endgroup

现在断言, \(F_0\backslash G^0\) 等同于 \(F\backslash\mathscr S(G^0)\).
证明与\SGARefLink{lem:sga3-V-6.1}{第~V~篇引理~6.1} 的断言~(i)
第一部分的证明类似. 考虑图
\begin{center}
  \phantomsection\label{diag:sga3-VIA-4.3-span}%
  \begin{tikzcd}[column sep=large]
    \mathscr S(G^0) &
    F\times_A G^0
      \arrow[l,"v"']
      \arrow[r,"\operatorname{pr}_2"] &
    G^0
  \end{tikzcd},
\end{center}
其中 \(v\) 为 \(\lambda\) 诱导的态射.
由于 \(\operatorname{pr}_2\) 有截面, 它是普遍有效满态射.
因此, \(F_0\backslash G^0\) 与 \(\operatorname{Coker}(v_0,v_1)\) 相同, 其中
\begin{center}
  \phantomsection\label{diag:sga3-VIA-4.3-V-groupoid}%
  \begin{tikzcd}[column sep=large]
    V_2
      \arrow[r,shift left=1.4ex,"v'_2"]
      \arrow[r,"v'_1"]
      \arrow[r,shift right=1.4ex,"v'_0"'] &
    V_1
      \arrow[r,shift left=0.8ex,"v_1"]
      \arrow[r,shift right=0.8ex,"v_0"'] &
    V=F\times_A G^0
  \end{tikzcd}
\end{center}
是群胚 \(G^0_*\) 在 \(\operatorname{pr}_2\) 下的逆像
(参见第~V~篇 \S~3.a), 因而也是 \(\mathscr S(G^0_*)\) 在复合态射
\begin{center}
  \phantomsection\label{diag:sga3-VIA-4.3-composite}%
  \begin{tikzcd}[column sep=large]
    F\times_A G^0
      \arrow[r,"\text{包含}"] &
    F\times_A\mathscr S(G^0)
      \arrow[r,"\operatorname{pr}_2"] &
    \mathscr S(G^0)
  \end{tikzcd}.
\end{center}
下的逆像.

类似地, 由于 \(v\) 忠实平坦且拟紧,
\(F\backslash\mathscr S(G^0)\) 与群胚 \(\mathscr S(G^0_*)\)
经基变换 \(v\) 得到的逆像的余核相同.
由第~V~篇 \S~3.c, 这个逆像同构于 \(V_*\).
因此, \(G^0_*\) 到 \(\mathscr S(G^0_*)\) 的典范包含诱导同构
\[
  F_0\backslash G^0
    \xrightarrow{\ \sim\ }
  F\backslash\mathscr S(G^0).
\]

最后, \(F\backslash\mathscr S(G^0)\) 的形成与有限局部自由基变换可交换,
因为 \(F_0\backslash G^0\) 的形成具有同一性质
(参见编者注~\ref{note:sga3-VIA-37} 及下面的 \SGARefLink{par:sga3-VIA-4.6}{4.6}).
